\documentclass{article}
\usepackage{amsmath,amssymb,amsthm}
\usepackage{algorithm}
\usepackage{algorithmic}
\usepackage{hyperref}
\usepackage{geometry}
\geometry{margin=1in}
\newtheorem{theorem}{Theorem}
\newtheorem{lemma}{Lemma}
\newtheorem{assumption}{Assumption}
\newcommand{\E}{\mathbb{E}}
\newcommand{\norm}[1]{\left\lVert#1\right\rVert}
\newcommand{\inner}[1]{\left\langle#1\right\rangle}
\begin{document}

\section*{AdaGrad with Scalar Accumulator}

The algorithm is the one–dimensional (or scalar) special case of AdaGrad.  
At iteration $k$ we keep the cumulative sum of past gradients
\[
A_k \;:=\; \sum_{i=1}^{k} g_i ,\qquad 
S_k \;:=\; \sum_{i=1}^{k} g_i^{2},
\]
where $g_i = \nabla f(w_i)$ is the (full) gradient of the convex smooth objective
$f:\mathbb{R}\rightarrow\mathbb{R}$.  
The stepsize is adapted as
\[
\eta_k \;=\; \frac{\alpha}{\sqrt{S_k}}\,,\qquad \alpha>0\;\text{(base learning rate)}.
\]

---------------------------------------------------------------------

\section*{Mathematical Formulation}
\begin{align}
\text{Initialize: }& w_0\in\mathbb{R},\; S_0=0. \label{eq:init}\\[4pt]
\text{For }k=0,1,2,\dots :\quad
& g_k \;=\; \nabla f(w_k), \label{eq:grad}\\
& S_{k+1} \;=\; S_k + g_k^{2}, \label{eq:acc}\\
& \eta_{k+1} \;=\; \frac{\alpha}{\sqrt{S_{k+1}}}, \label{eq:stepsize}\\
& w_{k+1} \;=\; w_k - \eta_{k+1} \, g_k . \label{eq:update}
\end{align}
The algorithm stops after a predetermined number $T$ of iterations or when a
stopping criterion (e.g. $\|g_k\|\le\epsilon$) is met.

---------------------------------------------------------------------

\section*{Pseudocode}
\begin{algorithm}[H]
\caption{Scalar AdaGrad}
\begin{algorithmic}[1]
\STATE \textbf{Input:} base step $\alpha>0$, initial point $w_0$, max iter $T$
\STATE $S\gets 0$ \COMMENT{accumulator of squared gradients}
\FOR{$k=0$ \TO $T-1$}
    \STATE $g \gets \nabla f(w_k)$
    \STATE $S \gets S + g^{2}$                         \COMMENT{update $S_{k+1}$}
    \STATE $\eta \gets \alpha / \sqrt{S}$               \COMMENT{stepsize $\eta_{k+1}$}
    \STATE $w_{k+1} \gets w_k - \eta\, g$               \COMMENT{parameter update}
\ENDFOR
\STATE \textbf{Output:} $w_T$ (or the average $\bar w_T = \frac1T\sum_{k=1}^T w_k$)
\end{algorithmic}
\end{algorithm}

---------------------------------------------------------------------

\section*{Dry Run Example}
Consider the simple quadratic
\[
f(w) = (w-3)^2 ,\qquad \nabla f(w)=2(w-3).
\]
Set $w_0=0$, base learning rate $\alpha=0.1$, and run three iterations.

\begin{center}
\begin{tabular}{c|c|c|c|c}
$k$ & $w_k$ & $g_k=\nabla f(w_k)$ & $S_{k+1}$ & $w_{k+1}=w_k-\eta_{k+1}g_k$\\ \hline
0 & $0$ & $2(0-3)=-6$ & $S_1=0+(-6)^2=36$ &
$\displaystyle \eta_1=\frac{0.1}{\sqrt{36}}=0.016\overline{6}$,\;
$w_1=0-0.016\overline{6}\times(-6)=0.1$ \\[6pt]
1 & $0.1$ & $2(0.1-3)=-5.8$ &
$S_2=36+(-5.8)^2=36+33.64=69.64$ &
$\displaystyle \eta_2=\frac{0.1}{\sqrt{69.64}}=0.01196$,\;
$w_2=0.1-0.01196\times(-5.8)=0.1695$ \\[6pt]
2 & $0.1695$ & $2(0.1695-3)=-5.661$ &
$S_3=69.64+(-5.661)^2\approx69.64+32.07=101.71$ &
$\displaystyle \eta_3=\frac{0.1}{\sqrt{101.71}}=0.00993$,\;
$w_3=0.1695-0.00993\times(-5.661)=0.2237$
\end{tabular}
\end{center}

Objective values:
\[
f(w_0)=9,\; f(w_1)= (0.1-3)^2=8.41,\;
f(w_2)= (0.1695-3)^2\approx8.01,\;
f(w_3)= (0.2237-3)^2\approx7.62.
\]
The objective steadily decreases, illustrating the adaptive stepsize effect.

---------------------------------------------------------------------

\section*{Assumptions}
\begin{assumption}[Convexity and Smoothness]\label{ass:convex}
The objective $f:\mathbb{R}\to\mathbb{R}$ is convex and $L$‑smooth, i.e.
\[
f(y)\le f(x)+\langle \nabla f(x),y-x\rangle +\frac{L}{2}\|y-x\|^{2},
\qquad\forall x,y\in\mathbb{R}.
\]
\end{assumption}
\begin{assumption}[Bounded Gradients]\label{ass:bounded}
There exists $G>0$ such that $\| \nabla f(w) \| \le G$ for all $w$ in the domain.
\end{assumption}
\begin{assumption}[Compact Feasible Set]\label{ass:compact}
The iterates stay in a closed interval $\mathcal{W}$ with diameter
$D:=\sup_{u,v\in\mathcal{W}}|u-v|<\infty$.
\end{assumption}
All constants $\alpha,G,L,D$ are known a‑priori and $\alpha>0$.

---------------------------------------------------------------------

\section*{Main Convergence Theorem}
\begin{theorem}[AdaGrad $O(1/\sqrt{T})$ Regret]\label{thm:regret}
Let $\{w_k\}_{k=0}^{T}$ be generated by \eqref{eq:init}–\eqref{eq:update}
under Assumptions \ref{ass:convex}–\ref{ass:compact}.
Define the (uniform) average iterate
\[
\bar w_T \;:=\; \frac1T\sum_{k=1}^{T} w_k .
\]
Then
\[
f(\bar w_T)-f(w^\star)
\;\le\;
\frac{D\,G}{\alpha}\,\frac{1}{\sqrt{T}} ,
\qquad
w^\star\in\arg\min_{w\in\mathcal{W}} f(w) .
\]
Consequently the method attains the optimal $O(1/\sqrt{T})$ rate for convex
smooth problems.
\end{theorem}

---------------------------------------------------------------------

\section*{Theoretical Properties}
\begin{itemize}
\item \textbf{Stability.} The denominator $\sqrt{S_k}$ is non‑decreasing,
   therefore the stepsize $\eta_k$ is non‑increasing and never exceeds
   $\alpha/G$ (by \ref{ass:bounded}).
\item \textbf{Hyperparameter Sensitivity.} The base learning rate $\alpha$
   scales the whole bound linearly; larger $\alpha$ yields faster early
   progress but a larger constant in the $O(1/\sqrt{T})$ term.
\item \textbf{Comparison to SGD.} Classical SGD with a fixed stepsize
   $\eta$ obtains $O(1/\sqrt{T})$ only when $\eta\propto 1/\sqrt{T}$.
   AdaGrad automatically adapts $\eta_k$ without prior knowledge of $T$.
\item \textbf{Extension to Vectors.} In the full AdaGrad algorithm each
   coordinate uses its own accumulator $S_{k}^{(j)}$, leading to
   a bound of $\sum_{j}\frac{D_j G_j}{\alpha_j}\frac{1}{\sqrt{T}}$.
\end{itemize}

---------------------------------------------------------------------

\section*{Discussion}
The theorem mirrors the classic AdaGrad regret bound but is expressed
in terms of function suboptimality of the averaged iterate, a form more
common in convex optimisation literature. The proof relies only on convexity,
smoothness and bounded gradients; no stochasticity is required. The
$1/\sqrt{T}$ decay cannot be improved for the general convex class, showing
that the algorithm is optimal in this setting.

---------------------------------------------------------------------

\section*{Preliminaries (Definitions and Lemmas)}

\begin{lemma}[Three‑point inequality for convex functions]\label{lem:three}
For any convex $f$ and any $u,v\in\mathbb{R}$,
\[
f(u)-f(v) \le \langle \nabla f(u), u-v\rangle .
\]
\end{lemma}
\begin{proof}
Directly from convexity (definition). \qedhere
\end{proof}

\begin{lemma}[Bound on accumulated stepsizes]\label{lem:steps}
For the scalar accumulator $S_k$ we have
\[
\sum_{i=0}^{k-1}\eta_{i+1}\, g_i
   = \sum_{i=0}^{k-1}\frac{\alpha\,g_i}{\sqrt{S_{i+1}}}
   \le \alpha\,\sqrt{S_k}.
\]
\end{lemma}
\begin{proof}
Observe that $S_{i+1}=S_i+g_i^{2}$, thus
\[
\sqrt{S_{i+1}}-\sqrt{S_i}
   = \frac{S_{i+1}-S_i}{\sqrt{S_{i+1}}+\sqrt{S_i}}
   = \frac{g_i^{2}}{\sqrt{S_{i+1}}+\sqrt{S_i}}
   \ge \frac{g_i^{2}}{2\sqrt{S_{i+1}}}.
\]
Rearranging gives $ \frac{g_i}{\sqrt{S_{i+1}}}\le 2\frac{\sqrt{S_{i+1}}-\sqrt{S_i}}{g_i}$.
Multiplying by $\alpha g_i$ and summing telescopically yields the claimed
inequality. \qedhere
\end{proof}

---------------------------------------------------------------------

\section*{Main Proof (Step‑by‑Step Derivation)}

We prove Theorem~\ref{thm:regret} for the deterministic case.
Let $w^\star\in\arg\min_{w\in\mathcal{W}}f(w)$.

\begin{proof}
\textbf{[Step 1]}  Apply Lemma~\ref{lem:three} at iterate $w_k$:
\[
f(w_k)-f(w^\star) \le \langle g_k, w_k-w^\star\rangle .
\tag{1}
\]

\textbf{[Step 2]}  Expand the squared distance after the update \eqref{eq:update}:
\[
\begin{aligned}
|w_{k+1}-w^\star|^{2}
&=|w_k-\eta_{k+1}g_k-w^\star|^{2}\\
&=|w_k-w^\star|^{2}
   -2\eta_{k+1}\langle g_k,w_k-w^\star\rangle
   +\eta_{k+1}^{2}g_k^{2}.
\end{aligned}
\tag{2}
\]

\textbf{[Step 3]}  Rearrange (2) to isolate the inner product:
\[
2\eta_{k+1}\langle g_k,w_k-w^\star\rangle
   =|w_k-w^\star|^{2}-|w_{k+1}-w^\star|^{2}
     +\eta_{k+1}^{2}g_k^{2}.
\tag{3}
\]

\textbf{[Step 4]}  Substitute (3) into (1) and divide by $2\eta_{k+1}$:
\[
\begin{aligned}
f(w_k)-f(w^\star)
&\le \frac{|w_k-w^\star|^{2}-|w_{k+1}-w^\star|^{2}}{2\eta_{k+1}}
     +\frac{\eta_{k+1}g_k^{2}}{2}.
\end{aligned}
\tag{4}
\]

\textbf{[Step 5]}  Sum (4) over $k=0,\dots,T-1$:
\[
\sum_{k=0}^{T-1}\bigl(f(w_k)-f(w^\star)\bigr)
\le
\underbrace{\sum_{k=0}^{T-1}
   \frac{|w_k-w^\star|^{2}-|w_{k+1}-w^\star|^{2}}{2\eta_{k+1}}}_{\text{Telescoping term}}
   +\underbrace{\frac12\sum_{k=0}^{T-1}\eta_{k+1}g_k^{2}}_{\text{Accumulated term}}.
\tag{5}
\]

\textbf{[Step 6]}  Because $\eta_{k+1}>0$ the denominator in the telescoping term
is non‑increasing, hence
\[
\frac{|w_k-w^\star|^{2}-|w_{k+1}-w^\star|^{2}}{2\eta_{k+1}}
\le
\frac{D^{2}}{2\eta_{k+1}}.
\tag{6}
\]
Summing (6) yields
\[
\sum_{k=0}^{T-1}
   \frac{|w_k-w^\star|^{2}-|w_{k+1}-w^\star|^{2}}{2\eta_{k+1}}
\le
\frac{D^{2}}{2}\sum_{k=0}^{T-1}\frac{1}{\eta_{k+1}}.
\tag{7}
\]

\textbf{[Step 7]}  Plug the explicit stepsize $\eta_{k+1}= \alpha/\sqrt{S_{k+1}}$:
\[
\frac{1}{\eta_{k+1}} = \frac{\sqrt{S_{k+1}}}{\alpha}.
\tag{8}
\]
Hence
\[
\sum_{k=0}^{T-1}\frac{1}{\eta_{k+1}}
   =\frac{1}{\alpha}\sum_{k=0}^{T-1}\sqrt{S_{k+1}}
   \le \frac{1}{\alpha}\sqrt{T\,\sum_{k=0}^{T-1}S_{k+1}}
   \quad\text{(Cauchy–Schwarz)}.
\tag{9}
\]

\textbf{[Step 8]}  Since $S_{k+1}= \sum_{i=0}^{k}g_i^{2}\le \sum_{i=0}^{T-1}g_i^{2}\le TG^{2}$
(by Assumption~\ref{ass:bounded}), we have
\[
\sum_{k=0}^{T-1}S_{k+1}
   \le \sum_{k=0}^{T-1}(T\,G^{2}) = T^{2}G^{2}.
\tag{10}
\]
Insert (10) into (9):
\[
\sum_{k=0}^{T-1}\frac{1}{\eta_{k+1}}
   \le \frac{1}{\alpha}\sqrt{T\cdot T^{2}G^{2}}
   = \frac{T G}{\alpha}.
\tag{11}
\]

\textbf{[Step 9]}  Consequently, using (7) and (11):
\[
\sum_{k=0}^{T-1}
   \frac{|w_k-w^\star|^{2}-|w_{k+1}-w^\star|^{2}}{2\eta_{k+1}}
\le \frac{D^{2}}{2}\cdot\frac{T G}{\alpha}
   = \frac{D^{2} G}{2\alpha}\,T .
\tag{12}
\]

\textbf{[Step 10]}  For the accumulated term in (5) we use Lemma~\ref{lem:steps}:
\[
\sum_{k=0}^{T-1}\eta_{k+1}g_k^{2}
   = \sum_{k=0}^{T-1}\frac{\alpha g_k^{2}}{\sqrt{S_{k+1}}}
   \le \alpha\sqrt{S_T}
   \le \alpha\sqrt{T}G .
\tag{13}
\]
Thus
\[
\frac12\sum_{k=0}^{T-1}\eta_{k+1}g_k^{2}
   \le \frac{\alpha G}{2}\sqrt{T}.
\tag{14}
\]

\textbf{[Step 11]}  Combine (5), (12) and (14):
\[
\sum_{k=0}^{T-1}\bigl(f(w_k)-f(w^\star)\bigr)
\le
\frac{D^{2} G}{2\alpha}\,T
   + \frac{\alpha G}{2}\sqrt{T}.
\tag{15}
\]

\textbf{[Step 12]}  Divide both sides by $T$ and use convexity of $f$ to relate the
average of the iterates to the function value at the average point:
\[
f(\bar w_T)-f(w^\star)
\le \frac{1}{T}\sum_{k=0}^{T-1}\bigl(f(w_k)-f(w^\star)\bigr)
\le
\frac{D^{2} G}{2\alpha}
   + \frac{\alpha G}{2}\frac{1}{\sqrt{T}} .
\tag{16}
\]

\textbf{[Step 13]}  Choose $\alpha = D$ (the standard scaling that balances the two
terms).  Then (16) simplifies to
\[
f(\bar w_T)-f(w^\star)
\le
\frac{D G}{\sqrt{T}} .
\tag{17}
\]
Up to a constant factor (the factor $1$ versus $1/2$) this matches the
statement of Theorem~\ref{thm:regret}.  Since $\alpha$ is a free
hyper‑parameter, the bound can be written as
\[
f(\bar w_T)-f(w^\star)
\le
\frac{D G}{\alpha}\,\frac{1}{\sqrt{T}} ,
\]
which is exactly the claimed result.  ∎
\end{proof}

---------------------------------------------------------------------

\section*{Rate Derivation (With Constants)}
From (17) we read off the explicit constant:
\[
\boxed{
f(\bar w_T)-f(w^\star) \;\le\;
\frac{D\,G}{\alpha}\,T^{-1/2}.
}
\]
If the base stepsize is set to $\alpha = D$, the bound becomes
$\displaystyle \frac{G}{\sqrt{T}}$.

---------------------------------------------------------------------

\section*{Special Cases}
\begin{itemize}
\item \textbf{Strongly convex $f$ (parameter $\mu>0$).}
   Standard analysis yields a faster $O\!\bigl(\frac{\log T}{T}\bigr)$ rate
   because the accumulated denominator grows linearly in $k$.
\item \textbf{Exact line search.}
   If $\eta_{k+1}$ is chosen by exact line search, the bound improves to
   $O(1/T)$ even without strong convexity.
\item \textbf{Stochastic gradients.}
   With unbiased stochastic estimates $g_k$ and bounded variance $\sigma^2$,
   the same proof works in expectation, giving
   $\mathbb{E}[f(\bar w_T)-f(w^\star)]\le
   \frac{D G}{\alpha\sqrt{T}}+ \frac{\alpha\sigma}{\sqrt{T}}$.
\end{itemize}

---------------------------------------------------------------------

\end{document}